\section{Balanced fibres beyond transitivity}
\label{sec:orbit-residues}

Transitivity in \cref{thm:balanced-pair} is a sufficient way to obtain
an orbit divisor.  It is not necessary.  The statements below use an
actual invariant scalar-word action and derive the tuple counts through
stabilizers.  They do not assume transitivity on tuples or independent
behaviour of the row and lane actions.

\subsection{The fixed-word residue test}

Let a group $U$ act on a finite scalar-word set $\Omega$, and let $I$
be a finite index set.  The action on $\Omega^I$ is diagonal.  Write
$h=|\operatorname{Fix}_U(\Omega)|$ and $r=|I|$.

\begin{theorem}[Fixed-word residues]
\label{thm:fixed-word-residue}
\formal{CompositionOrbitResidues.lean}{tuple_orbit_divisor,fixedTupleEquiv,invariant_fiber_fixed_residue,balanced_fixed_word_residue}
Suppose $m$ divides the cardinality of every non-singleton $U$-orbit
in $\Omega$.  Let $f\colon\Omega^I\to Z$ be $U$-invariant, where
$Z$ is finite, and suppose every fibre has the same cardinality $F$.
Then
\[
 |Z|\,(F\bmod m)\le h^r.
\]
In particular, in the complete matrix interface over a scalar alphabet
of size $a$, with $r=db$, $|Z|=a^{db}$ and
$F=a^{db(d-1)}$, any such invariant action satisfies
\begin{equation}\label{eq:fixed-word-residue}
 h^{db}\ \ge\ a^{db}\bigl(a^{db(d-1)}\bmod m\bigr).
\end{equation}
The integer $m$ need not be prime.
\end{theorem}
\begin{outline}
If a tuple $x$ is not fixed, some coordinate $x_i$ is not fixed.
The stabilizer containment $U_x\le U_{x_i}$ shows that
$|U\cdot x_i|$ divides $|U\cdot x|$.  Therefore every non-singleton
tuple orbit has size divisible by $m$.

An invariant fibre is a union of tuple orbits.  Its cardinality modulo
$m$ equals the number of its fixed tuples modulo $m$.  Consequently
each fibre contains at least $F\bmod m$ fixed tuples.  Fixed tuples
are precisely tuples of fixed scalar words, of which there are $h^r$.
Summing over the disjoint fibres proves the inequality.  In the complete
matrix setting, bijectivity gives the stated fibre size, as in the proof
of \cref{thm:balanced-pair}.
The formal argument proves the class equation for an arbitrary common
orbit divisor and constructs the equivalence between fixed tuples and
tuples of fixed words.
\end{outline}

\begin{corollary}[Too few fixed words]
\label{cor:few-fixed-words}
\formal{CompositionOrbitResidues.lean}{balanced_fixed_word_small_impossible}
Under the preceding hypotheses, suppose $r>0$, $m>1$,
$\gcd(m,a)=1$, $|Z|=a^r$, and $F=a^n$ for some $n\ge0$.
If $h<a$, no such invariant balanced function exists.
In particular this applies to any prime divisor $m$ coprime to $a$.
\end{corollary}
\begin{outline}
Coprimality implies $m\nmid a^n$, so $a^n\bmod m\ge1$.
The residue inequality gives $a^r\le h^r$, whereas $h<a$ and
$r>0$ give the strict reverse inequality.
\end{outline}

There is a useful version requiring only one fibre, without counting
the fixed set.

\begin{proposition}[A semiregular divisor]
\label{prop:semiregular-balance}
\formal{CompositionOrbitResidues.lean}{no_balanced_semiregular,no_balanced_free_off_zero}
Suppose $U$ fixes $0\in\Omega$ and every orbit of a word other than
zero has cardinality divisible by $m>1$, with $\gcd(m,a)=1$.
If $f\colon\Omega^I\to Z$ is invariant, no fibre avoiding the
all-zero tuple has cardinality $a^n$, for any $n\ge0$.
In particular, a cyclic action of order $m$ that is free away from zero
has this obstruction.  Hence it cannot be an invariance subgroup of a
nonconstant balanced decoder whose fibre sizes are powers of $a$.
\end{proposition}
\begin{outline}
Every tuple in the selected fibre has a nonzero coordinate.  Its
stabilizer is contained in that coordinate's stabilizer, so its orbit
size is a multiple of $m$.  The invariant fibre is a disjoint union of
these orbits and also has size divisible by $m$.  This contradicts
$\gcd(m,a^n)=1$.  For a free action of a finite cyclic group of order
$m$, each nonzero scalar orbit has size exactly $m$, giving the stated
specialization.  The proof uses this orbit-size condition directly; it
does not assume that all nonzero scalar words belong to one orbit.
\end{outline}

\subsection{An actual nontransitive order-five example}

Set $d=4$ and $\Omega=\F_2^4$, with cyclic shift
$(Tx)_i=x_{i+1}$.  Define the literal word operation
\[
 L(x_0,x_1,x_2,x_3)
   =(x_2,\ x_2+x_3,\ x_1+x_2,\ x_0+x_2),
 \qquad Q_c=T^cLTL^{-1}T^{-c}.
\]
All equalities in the next result concern these displayed operations;
no assumed scalar-group identification is used.

\begin{theorem}[Balance excludes a nontransitive adjacent pair]
\label{thm:order-five-example}
\formal{CompositionOrderFive.lean}{orderFivePerm_zero,orderFivePerm_fixed,orderFivePerm_order,orderFive_conjugate_difference,orderFive_pair_nontransitive,orderFive_no_balanced_covariance,orderFive_adjacent_class_impossible}
The operation $L$ is a zero-fixing permutation of order five, and fixes
no other word.  The group $\langle Q_0,Q_1\rangle$ is not transitive
on the fifteen nonzero words.
Nevertheless, for every $b\ge1$, a complete factorization
$KLH=\tau$ on $\F_2^{4\times4\times b}$, with $H$ equivariant and
$K$ class-uniform, must satisfy $p(0)\ne p(1)$.
The same scalar $L$ is applied independently to every row and lane;
$H$ and $K$ may be nonlinear and may mix all lanes.
\end{theorem}
\begin{outline}
Substitution gives
\[
 L^4x=(x_0+x_3,\ x_0+x_2,\ x_0,\ x_0+x_1),
 \qquad L(L^4x)=L^4(Lx)=x.
\]
The fixed-point equations give successively $x_0=x_2$, $x_1=0$,
$x_3=0$, and $x_2=0$.  Thus $L$ fixes only zero, and since
$L^5=I$ and $L\ne I$, its order is five.  The actual conjugates are
\[
 \begin{aligned}
 Q_0x&=(x_0+x_1,\ x_1+x_3,\ x_1,\ x_1+x_2),\\
 Q_1x&=(x_0+x_2,\ x_0,\ x_0+x_1,\ x_0+x_3).
 \end{aligned}
\]
They satisfy
\[
 Q_0Q_1^{-1}=L.
\]
Write $e_i$ for the coordinate unit words and $v=(1,1,1,1)$.
On $\{e_0,e_1,e_2,e_3,v\}$, the first conjugate fixes $e_0$
and cycles $(e_1,v,e_2,e_3)$; the second fixes $e_3$ and cycles
$(e_0,v,e_1,e_2)$.  Thus both preserve this five-word set.  In particular no group element sends
a unit word to a word with exactly two nonzero coordinates, so the pair
group is not transitive on the nonzero words.  Lean checks these finite
identities on the sixteen literal words in its kernel.

If $p(0)=p(1)$, the common shared decoder $f$ has
$f(Q_0Y)=\sigma f(Y)=f(Q_1Y)$.  Substitute $Q_1^{-1}Y$ to obtain
$f(LY)=f(Y)$, and hence invariance under the cyclic group generated by
$L$.  Because $L$ has prime order five and fixes only zero, every
nonzero scalar orbit has size five.  Stabilizer containment therefore
makes every nonzero tuple orbit size divisible by five.

Exactness on the complete finite space makes $H$ and $K$ bijective.
The column decoder is balanced: each value has exactly
$2^{12b}$ preimages in the $4b$-tuple domain.
Choose an output value different from $f(0)$; it exists because the
column alphabet has $2^{4b}>1$ values.  Its fibre avoids zero and is
invariant under $L$, so the semiregular-divisor proposition says its
size is divisible by five.  The power $2^{12b}$ is not, a contradiction.
The formal endpoint derives covariance and balance from the full product
identity; neither property is supplied as a separate certificate.
\end{outline}

These tests are necessary conditions for an actual complete factorization.
Passing an orbit-assignment test for one shared function does not produce
jointly independent output columns or a bijective whole-matrix map.
The arguments also impose no cost bound on evaluating a scalar-word
permutation or an orbit representative.

\subsection{Field multipliers and two scope checks}

\begin{proposition}[Multiplier class costs]\label{prop:multiplier-classes}
\formal{CompositionMultiplierClasses.lean}{fieldMultiplierConjugate_eq,multiplier_class_residue,multiplier_class_count,multiplier_class_count_prefix}
\formal{CompositionMultiplierGetters.lean}{actual_multiplier_getters_card,multiplier_degree_le,multiplier_binary_degree_eq}
In the balanced-intertwiner interface, identify scalar words over $\F_q$
with $E=\F_{q^d}$ in a reversed normal basis, so rotation is
$\Phi(x)=x^q$. Let $L(x)=\lambda x$, $\lambda\ne0$, on each row and lane.
Put $\kappa=\lambda^{1-q}$, $h=[\F_q(\kappa):\F_q]$ and
$h_0=[\F_q(\lambda):\F_q]$. Every second-layer class is contained in
one residue class modulo $h$, so it requires at least $h$ classes.
More generally, among the first $n$ receivers there are at least
$\min(n,h)$ distinct class labels, for every $0\le n\le d$.
The complete operation $L$ itself requires exactly $h_0$ distinct class
getters, and $h\le h_0$. For $q=2$, $h=h_0$.
\end{proposition}
\begin{outline}
Writing $m_z$ for scalar multiplication, direct composition gives
\[
 Q_c=m_{\Phi^c(\kappa)}\Phi,\qquad
 Q_cQ_e^{-1}=m_{\Phi^c(\kappa)/\Phi^e(\kappa)}.
\]
If $c,e$ share a decoder, its covariance makes it invariant under the
second map. A multiplier $\mu\ne1$ has order $t>1$ dividing $q^d-1$;
its cyclic action is free away from zero. Since $\gcd(t,q)=1$,
\cref{prop:semiregular-balance} contradicts the required balance. Thus
$\Phi^c(\kappa)=\Phi^e(\kappa)$, equivalent to $c\equiv e\pmod h$.
The first $\min(n,h)$ receivers have distinct residues, proving the prefix
bound without requiring a full extension-degree set of receivers.

Let $\ell_0:E\to\F_q$ extract coordinate zero. The unique full-domain
getter for $L$ at receiver $c$ is, on each row and lane,
$x\mapsto\ell_0(\Phi^c(\lambda)x)$: rotate the received word back,
apply $L$, then rotate its demanded coordinate to zero. Two such maps
agree exactly when their scalar multipliers agree, since a nonzero
scalar difference is invertible and $\ell_0\ne0$. There are $h_0$
conjugates, giving both necessity and sufficiency. The inclusion
$\kappa\in\F_q(\lambda)$ gives $h\le h_0$; for $q=2$,
$\kappa=\lambda^{-1}$ gives equality. Thus this family does not hide a
new small-first-layer exclusion. For odd $q$, a nonzero trace-zero
$\lambda\in\F_{q^2}$ has $h_0=2$ but $\kappa=-1$ and $h=1$,
so the distinction is necessary.
\end{outline}

\begin{example}[Surjectivity is weaker than balance]\label{ex:ratio-unbalanced}
\formal{CompositionRatio.lean}{ratioDecoder_surjective,ratio_nonzero_fiber_card,ratio_zero_fiber_card,ratio_eight_fiber_cards}
\formal{CompositionMultiplierClasses.lean}{ratio_multiplier_conjugate_covariance}
For $E=\F_8$, $d=3$, $b=1$ and any $\lambda\ne0$, use the preceding
actual conjugates. On the complete row-word input $E^3$, the total map
\[
 f(x_0,x_1,x_2)=
 \begin{cases}x_1/x_0&x_0\ne0,\\0&x_0=0\end{cases}
\]
is surjective and satisfies $f(Q_cX)=\Phi(f(X))$ for every receiver.
Its seven nonzero fibres have size $56$ and its zero fibre has size $120$,
rather than the required common size $64$. More generally, over a field
of size $Q$, the corresponding fibre sizes are $(Q-1)Q$ and
$Q^2+(Q-1)Q$.
\end{example}
\begin{outline}
Each $Q_c$ applies Frobenius and the same nonzero multiplier to all three
rows, so the common multiplier cancels in the quotient; the zero branch
is preserved. The input $(1,z,0)$ attains any $z$. For nonzero $z$ there
are seven choices of $x_0$ and eight of $x_2$, with $x_1=zx_0$ fixed.
For $z=0$, add the $64$ inputs with $x_0=0$ to the $56$ with
$x_0\ne0,x_1=0$. These sum to $512$. No complete bijective transposer
or balanced decoder is constructed by this example.
\end{outline}

\begin{proposition}[Exact orbit-assignment criterion]\label{prop:orbit-assignment}
\formal{CompositionOrbitAssignment.lean}{OrbitAssignment.orbit_map_exists_iff,OrbitAssignment.allowed_iff_divides,OrbitAssignment.cyclic_balanced_map_iff_orbit_assignment}
Let a finite group $G$ act on a finite source set and let
$\theta:G\twoheadrightarrow\mathbb Z/d$ specify the cyclic action on a
finite target set. For each source orbit $O$, choose $x\in O$ and put
$n_O=|O|$ and $t_O=[\mathbb Z/d:\theta(G_x)]$.
For each individual target orbit $D$, put $\ell_D=|D|$.
An equivariant map with every target fibre of size $F$ exists if and only
if the whole source orbits can be assigned to individual target orbits,
with $\ell_D\mid t_O$ for every assignment and
\[
 \sum_{O\ \mathrm{assigned\ to}\ D}\frac{n_O}{\ell_D}=F
 \quad\text{for each }D.
\]
\end{proposition}
\begin{outline}
A map from the transitive orbit $O$ to $D$ is determined by the image $z$
of $x$. It is well defined precisely when
$\theta(G_x)\subseteq(\mathbb Z/d)_z$, equivalently
$\ell_D\mid t_O$ by the subgroup structure of a finite cyclic group.
When it exists, the map $gx\mapsto\theta(g)z$ is onto $D$ and every
fibre has size $n_O/\ell_D$: group translation bijects any two fibres,
and counting their disjoint union gives the quotient. The phase
stabilizer is independent of the source representative because the
phase group is abelian. Restricting a balanced map to each source
orbit proves necessity. Conversely choose one such orbit map for every
assignment; their disjoint union has the stipulated fibre sizes.
Whole source orbits cannot be split fractionally. This is an exact test
for one class function, not for joint bijectivity of all output columns.
\end{outline}
